\documentclass[../infufrgs/ppgc-tese.tex]{subfiles}
\begin{document}

\begin{flushright}
\mbox{}%\vfill
{\sffamily\itshape
``No matter how correct a mathematical theorem may appear to be, one ought never to be satisfied that there was not something imperfect about it until it also gives the impression of being beautiful.'' \\}
--- \textsc{George~Boole}
\end{flushright}

\section{LTN Cross-layer Transcoder}
The methodology combines three components: sparse feature extraction, cross-layer alignment, and logic-based validation/steering. The design is grounded on two completed lines of work: (i) controlled late-generalization experiments with distribution shifts (ICTAI paper), and (ii) sparse-autoencoder analyses showing that small latent subsets can separate clean and perturbed samples in vision models. In this thesis, we transfer these insights to attention-based models under a unified pipeline.

Given a model and an input sequence, we collect residual activations $h_\ell$ at selected layers. We train sparse autoencoders (SAEs) to obtain sparse codes $s_\ell$, then learn cross-layer transcoders $T_{\ell\to k}$ to map features across depth. On top of $s_\ell$ and $\hat s_k$, we define soft logic validators (LTN-style predicates/rules) that produce differentiable satisfaction scores. These scores are used both for analysis (semantic consistency) and for control signals during inference (e.g., reranking, constrained decoding, and logit adjustments).

The key methodological claim is that local behavior is better diagnosed when feature-level geometry, cross-layer consistency, and rule satisfaction are measured together, rather than separately.

\section{Implementation Strategy}
\textbf{Stage 1: Reproducible data and model setup.}
We instantiate three families of experiments: (a) synthetic class/subclass datasets that induce measurable distribution shifts and grokking-like transitions; (b) language-generation tasks for steering; and (c) robustness tasks with adversarial/corruption perturbations. This follows the principle validated in previous work: stress-test the mechanism in controlled environments first, then transfer to more realistic benchmarks.

\textbf{Stage 2: Layerwise sparse feature learning.}
For each selected layer $\ell$, we train an SAE with reconstruction and sparsity objectives. We track reconstruction error, activation sparsity, dead-feature rate, and feature stability across seeds/epochs. This stage provides the canonical latent representation used downstream.

\textbf{Stage 3: Cross-layer alignment.}
We train $T_{\ell\to k}$ to predict layer-$k$ sparse codes from layer-$\ell$ sparse codes. Optimization combines alignment loss and regularization, with optional validator consistency penalties. We evaluate whether mappings preserve both numeric alignment and semantic behavior under validators.

\textbf{Stage 4: Logic-guided evaluation and steering.}
We define LTN-style predicates/rules on sparse features (and optionally outputs), producing satisfaction scores in $[0,1]$. These scores are used in two modes: evaluation-time diagnostics and inference-time steering. Ablations remove one component at a time (no SAE, no transcoder, no logic constraints) to isolate causal contribution.

\textbf{Stage 5: Statistical analysis and reproducibility.}
All reported results are aggregated over multiple seeds. We report confidence intervals for key metrics and publish scripts/configurations for exact reruns.

\section{Evaluation Tasks}
\textbf{Goal G1: Emergence and late generalization.}
\begin{itemize}
    \item Tasks: synthetic class/subclass benchmarks and grokking-prone setups from previous work.
    \item Behavioral metrics: train/test accuracy dynamics, delay to generalization, stability of post-transition performance.
    \item Mechanistic metrics: SAE sparsity/reconstruction, cross-seed feature stability, cross-layer alignment quality.
\end{itemize}

\textbf{Goal G2: Steering locality and trade-offs.}
\begin{itemize}
    \item Tasks: constrained generation with single-rule and multi-rule settings.
    \item Behavioral metrics: constraint adherence, task quality, regression on non-targeted behaviors.
    \item Mechanistic metrics: validator satisfaction, layerwise activation shifts, consistency of targeted vs collateral feature changes.
\end{itemize}

\textbf{Goal G3: Shift/attack diagnostics.}
\begin{itemize}
    \item Tasks: adversarial/corruption perturbation suites and controlled watermark-like perturbations.
    \item Behavioral metrics: detection ROC-AUC/PR-AUC, calibration, early-warning lead time.
    \item Mechanistic metrics: proportion of attack-sensitive sparse features, separability of clean vs perturbed latent codes, cross-layer signature persistence.
\end{itemize}

\textbf{Baselines and ablations.}
Baselines include output-only heuristics (logit margin and confidence), dense-activation detectors, and steering without logic constraints. Core ablations remove each pipeline module and vary sparsity targets to quantify robustness of conclusions.

\section{Observed Effects}
The completed work already supports feasibility of the methodology:
\begin{itemize}
    \item Distribution shifts can reliably induce late generalization in controlled datasets, supporting G1 experimental design.
    \item Sparse latent representations in vision models show clear clean-vs-perturbed separation with interpretable features, supporting G3 diagnostics design.
    \item Detection with a compact subset of sparse features can remain competitive with dense baselines, suggesting practical value of sparse local representations.
\end{itemize}

These observations motivate the final thesis experiments: extending sparse-feature diagnostics from vision and synthetic settings to attention-based steering and cross-layer logical validation in a single framework.
\end{document}
